% Agent 13: Section 4.2 Adiabatic Hydrodynamic Accretion (continued)
% Book pages 390--395 (PDF pages 24--27)
% Covers: Bernoulli equation (4.2.3), solution plane analysis, sonic point,
%         bisectrix, Figures 4.2.1 and 4.2.2, sonic radius, accretion rate
% Equations (4.2.3) through (4.2.15b)

% ---------------------------------------------------------------
%  This file is meant to be \input into the master document.
%  It continues Section 4.2 beginning at the Bernoulli equation.
%  The continuity equation (4.2.1) and Euler equation (4.2.2)
%  have already been introduced on the preceding pages.
% ---------------------------------------------------------------

[Cf.\ eqs.\ (2.5.23) and (2.5.59).]

Here $M_0$ is the total rate of accretion of rest mass, $u$ is the radial
velocity, and $M$ is the mass of the hole.  We assume the gas has constant
adiabatic index~$\Gamma$, so that during the accretion the pressure and
density are related by the adiabatic law $p = K\rho^\Gamma$, and the speed of
sound is given by $a = (\Gamma p)^{1/2}$.  Let the constants $K$ and $\Gamma$
be given.  Our task is to determine the accretion rate~$\dot{M}_0$ and to
compute as well the distributions of density $\rho(r)$ and velocity $u(r)$
in the flow.

In place of the Euler equation~(4.2.2) we shall use the Bernoulli
equation\footnote{Notice that our notation differs from that in \S\,2.
Throughout \S\,4 we denote (for ease of eyesight)
$a$ = (speed of sound) = $a$, not $c_s$;
$\rho$ = (density of rest mass) = $\rho$, not $\rho_0$.
No confusion is likely, since nowhere in \S\,4 shall we deal with
4-accelerations or with total mass-energy $\rho = \rho_0(1+\pi)$, and
because $\Gamma \ll 1$ in all regions of the accreting gas.
We shall also retain factors of $G$, $c$, and $k$ (i.e.\ use ``cgs'' units)
throughout \S\,4.}
(2.5.38) rewritten in the form
\begin{equation}
\tfrac{1}{2}\,u^2 + \frac{1}{\Gamma - 1}\,a^2 - \frac{GM}{r}
= \text{constant}
= \frac{1}{\Gamma - 1}\,a_\infty^2\,.
\tag{4.2.3}
\end{equation}

[Here the enthalpy has been written in the form
$w = \int \rho^{-1}\,dp = a^2/(\Gamma-1)$;
cf.\ eq.\ (2.5.42).]  The constant has been determined by conditions at
infinity, where the gas is at rest.

Rewrite the law of mass conservation~(4.2.1) with density expressed in terms
of sound speed
\begin{equation}
u = \frac{\dot{M}_0}{4\pi\rho_\infty\,r^2}
    \left(\frac{a_\infty}{a}\right)^{2/(\Gamma-1)}\!.
\tag{4.2.4}
\end{equation}

The only unknown parameter in the coupled equations~(4.2.3) and~(4.2.4) is
the accretion rate~$\dot{M}_0$.  Once $\dot{M}_0$ is known, one can readily
solve for the radial distributions of velocity, sound speed, and density.

The mass flux is determined in the following way.  Consider the system of
equations~(4.2.3) and~(4.2.4).  In the $u$--$a$ plane, the Bernoulli
equation~(4.2.3) for fixed radius~$r$ defines an ellipse; each value of $r$
corresponds to a different ellipse.  Similarly, the equation of mass
conservation~(4.2.4) for each value of $r$ defines a hyperbola of fractional
power, $u\,a^{2/(\Gamma-1)} = \text{const}$.  Now, let the accretion
rate~$\dot{M}_0$ be chosen arbitrarily.  For every value of $r$,
equations~(4.2.3) and~(4.2.4) are 2 equations with 2 unknowns, $u$ and $a$.
Solving these equations---i.e.\ finding the intersection point of the ellipse
with the hyperbola---gives $u$ and $a$ for that particular radius.  In other
words, in the $u$, $a$ plane the curve~$u(a)$ is determined parametrically by
the intersections of corresponding ellipses and hyperbolae.  It is clear (see
Figure~\ref{fig:4.2.1}) that for every ellipse-hyperbola pair, there are
either two intersection points, or one point of tangency, or no intersection
at all.  If, for a particular chosen~$\dot{M}_0$, there exists any $r$ at
which the two equations are incompatible for the~$\dot{M}_0$, such flow
cannot occur.  Rejecting this case, we have 2 remaining possibilities, for a
given~$\dot{M}_0$ (see Figure~\ref{fig:4.2.2}):

\begin{enumerate}
\item The ellipse and hyperbola intersect twice for every value of~$r$
  (Figure~\ref{fig:4.2.2}a).\\
  In this case we have 2 separate curves, corresponding to 2 families of
  intersection points.

\item The ellipse and hyperbola meet tangentially
  (Figure~\ref{fig:4.2.2}b) at one section point (call it~$r_s$), at which
  they cross each other; the two families of intersection points $u(a)$
  cross each other at $r = r_s$.  We shall call the two families of
  intersection points on the ``bisectrix'' of the graph, below that this
  crossing point necessarily lies on the line $u = a$.
\end{enumerate}

In both the cases, 1 and 2, the curves $u(a)$ describe possible flows of gas
in the gravitational field.  The curves begin on the innermost ellipse,
$r = \infty$.  On the upper curve~$u(a)$, at this ellipse we have $a = 0$,
but $u \neq 0$.  These boundary conditions are not of interest for the
accretion problem,\footnote{This curve describes the ``stellar wind'' by
which some stars eject matter into interstellar space; see Chapter~13 of
ZN.} because accretion requires $u_\infty = 0$,
$a_\infty \neq 0$.  The lower curve~$u(a)$ begins at the point
$u_\infty = 0$, $a_\infty \neq 0$---

\begin{figure}[htbp]
\centering
% \includegraphics[width=\columnwidth]{fig_4_2_1}
\fbox{\parbox{0.9\columnwidth}{\centering [Figure 4.2.1 placeholder]\\[6pt]
The $u$, $a$ ``solution plane'' for solving the coupled Bernoulli
equation~(4.2.3) and law of mass conservation~(4.2.4).
Solid curves: ellipses from the Bernoulli equation at several fixed values
of~$r$ ($r = \infty$ is the innermost ellipse).
Dashed curves: hyperbolae from the continuity equation.
Heavy dots mark intersection points; the locus of these intersections
traces the flow solution $u(a)$.
The diagonal $u = a$ is the ``bisectrix'' $D = n$.
Regions above the bisectrix are supersonic; regions below are subsonic.}}
\caption{The $u$, $a$ ``solution plane'' for solving the coupled Bernoulli
equation~(4.2.3) and law of mass conservation~(4.2.4).}
\label{fig:4.2.1}
\end{figure}

to our desired boundary condition.  Hence we shall restrict attention to the
lower curve.

We must still decide which type of lower curve is reasonable---one that
remains always below the bisectrix (case~1), or one that crosses the
bisectrix (case~2).  In case~1 the flow is subsonic at all radii, $u < a$.
But this requires a large back-pressure at all radii to retard the
inflow---back pressure that can never be provided near the gravitational
radius of a black hole.  The gas must cross the gravitational

\begin{figure}[htbp]
\centering
% \includegraphics[width=\columnwidth]{fig_4_2_2}
\fbox{\parbox{0.9\columnwidth}{\centering [Figure 4.2.2 placeholder]\\[6pt]
Possible forms of adiabatic gas flow in a spherical, Newtonian gravitational
field.  Curves~(a) [case~1 in text] correspond to flow that remains subsonic
or supersonic everywhere.  Curves~(b) [case~2 in text] correspond to flow
for which there is a sonic point, $u = a$.}}
\caption{Possible forms of adiabatic gas flow in a spherical, Newtonian
gravitational field.  Curves~(a) [case~1 in text] correspond to flow that
remains subsonic or supersonic everywhere.  Curves~(b) [case~2 in text]
correspond to flow for which there is a sonic point, $u = a$.}
\label{fig:4.2.2}
\end{figure}

\noindent
radius with the speed of light, as measured in the \emph{proper} reference
frame of an observer there; hence, the flow is surely supersonic near the
gravitational radius.

Thus, the only acceptable solution $u(a)$ is the lower curve of case~2
(Figure~\ref{fig:4.2.2}b).  This curve begins at $u = u_\infty = 0$,
$a = a_\infty \neq 0$, and it crosses from the subsonic region into the
supersonic region at $r = r_s$.  This transition to supersonic flow is
crucial to the accretion process.  It governs the accretion rate~$\dot{M}_0$,
in the same manner as the transition to supersonic flow in the throat of a
rocket nozzle.  In both cases one and only one mass flow rate is compatible
with the required transition to supersonic flow.

Let us calculate the required accretion rate~$\dot{M}_0$.  We begin by
calculating the flow velocity at the transition radius (``sonic
radius'')~$r_s$.  To do this, we rewrite the Euler equation~(4.2.2) in the
form
\begin{equation}
\frac{du}{dr} = -\,\frac{a^2\,dp}{p\,dr} - \frac{GM}{r^2}\,.
\tag{4.2.5}
\end{equation}

We differentiate the law of mass conservation~(4.2.1) with respect to $r$
and put $du/dr$ from it into~(4.2.5).  As a result we obtain
\begin{equation}
\frac{dp}{dr}\;\frac{(u^2 - a^2)}{p}
  = -\frac{GM}{r^2} + \frac{2u^2}{r}\,,
\tag{4.2.6}
\end{equation}
%
which shows that at the sonic point, where $a = u$,
\begin{equation}
2u_s^2 = \frac{GM}{r_s}\,,\qquad \text{or} \quad
u_s^2 = a_s^2 = \tfrac{1}{2}\,\frac{GM}{r_s}\,.
\tag{4.2.7}
\end{equation}

By combining this result with the Bernoulli equation~(4.2.3) we obtain the
speed of sound at the sonic point in terms of the speed of sound at infinity
\begin{equation}
a_s = u_s = a_\infty \left(\frac{2}{5-3\Gamma}\right)^{1/2}\!;
\tag{4.2.8}
\end{equation}
%
and using equation~(4.2.7) we obtain the radius at the sonic point
\begin{equation}
r_s = \left(\frac{5-3\Gamma}{4}\right)\frac{GM}{a_\infty^2}\,.
\tag{4.2.9}
\end{equation}
%
(Notice that the gravitational potential $GM/r_s$ at the sonic point is of
order~$a_\infty^2$.)

Now it is straightforward to calculate the accretion rate from
equation~(4.2.4):
\begin{equation}
\dot{M}_0 = 4\pi\,\alpha\,G^2 M^2\,\rho_\infty\,a_\infty^{-3}
= 4\Gamma^{3/2}\,\alpha\,G^2\,M^2\,\rho_\infty\,
  (m_p / k T_\infty)^{3/2}\,.
\tag{4.2.10}
\end{equation}

Here $\alpha$ is a constant of order unity which depends on~$\Gamma$:
\begin{equation}
\alpha \equiv \frac{\pi}{4\Gamma^{3/2}}
  \left(\frac{2}{5-3\Gamma}\right)^{(5-3\Gamma)/[2(\Gamma-1)]}\!.
\tag{4.2.11}
\end{equation}
%
\begin{align}
\alpha &\simeq 1.5  \quad \text{for} \quad \Gamma = 1,    \notag \\
\alpha &\simeq 1.2  \quad \text{for} \quad \Gamma = 1.4,  \notag \\
\alpha &\simeq 0.3  \quad \text{for} \quad \Gamma = 5/3;
\end{align}
%
and we have used the equation of state
$p_\infty = (\rho_\infty / m_p)\,k\,T_\infty$.

This expression for the accretion rate is the main result of our calculation.

Notice that $\dot{M}_0$ depends on~$\Gamma$.  At a typical temperature of
$T_\infty \sim 10^4$\,K the interstellar gas will be partially ionized, so
when it is compressed some energy will go into further ionization.  This means
that the value of $\Gamma$ outside and near the sonic point will be somewhat
below $5/3$.  A value of $\Gamma = 1.4$ might be reasonable for use in
equations~(4.2.10) and~(4.2.11).

It is now straightforward to derive from equations~(4.2.3) and~(4.2.4) the
radial distributions of all quantities.  The general picture is as follows.
Outside the ``radius of influence''
\begin{equation}
r_i = 2\,GM/a_\infty^2 = r_g\,(c/a_\infty)^2
\tag{4.2.12}
\end{equation}
%
\[
\simeq 10\,r_g \quad \text{for} \quad \Gamma = 1.4,
\]
%
the pull of the hole hardly makes itself felt, so $\rho$ and $a$ are nearly
equal to their values at ``infinity''; $\rho \simeq \rho_\infty$,
$a \simeq a_\infty$; and the density and sound velocity begin to rise.
After the gas passes through the sonic point~$r_s$, it is in near free-fall

\begin{equation}
u = (2GM/r)^{1/2} = a_\infty\,(r_i/r)^{1/2}
\qquad \text{at} \quad r < r_s\,;
\tag{4.2.13a}
\end{equation}
%
so the law of mass conservation requires the density to increase as
\begin{equation}
\rho = \frac{\dot{M}_0}{4\pi r^2\,u}
\simeq 0.2\,\rho_\infty\,(r_i/r)^{3/2}
\qquad \text{at} \quad r < r_s\,.
\tag{4.2.13b}
\end{equation}

For adiabatic compression, temperature rises as
$T \propto \rho^{\,\Gamma-1}$; consequently, at $r < r_s$
\begin{equation}
T = \left(\frac{a\Gamma^{3/2}}{4\pi}\right)^{(\Gamma-1)}
\simeq 0.5\,T_\infty\left(\frac{r_i}{r}\right)^{\frac{1}{2}(\Gamma-1)}
\qquad \text{if} \quad \Gamma = 1.4\;\text{near sonic point.}
\tag{4.2.13c}
\end{equation}

The above solution, (4.2.10)--(4.2.13), for adiabatic hydrodynamic accretion
is due originally to \citet{Bondi1952}.

Rewrite the accretion rate~(4.2.10) for hydrodynamic flow in a form similar
to that (eq.~4.1.3) for noninteracting particles
\begin{equation}
\dot{M}_0 = \pi\,r_g^2\,\rho_\infty\,c\,(c/a_\infty)^3.
\tag{4.2.14}
\end{equation}

Direct comparison, and use of the approximate equality between speed of sound
$a_\infty$ and proton speeds~$v_\infty$, shows that the hydrodynamic
accretion rate is larger by a factor $(c/a_\infty)^2 \simeq 10^{10}$ than
the accretion rate for independent particles.  The physical reason for this
is clear: gas is distinguished from independent particles by the frequent
collisions of its atoms; these collisions limit the tangential velocities
during infall, but permit radial velocities to grow.

Other, useful forms for the hydrodynamic accretion rate are
\begin{equation}
\dot{M}_0 \simeq 10^{5}
\left(\frac{M}{M_\odot}\right)^{\!2}
\left(\frac{\rho_\infty}{10^{-24}\;\text{g\,cm}^{-3}}\right)
\left(\frac{a_\infty}{10\;\text{km\,sec}^{-1}}\right)^{\!-3}
\;\text{,}
\tag{4.2.15a}
\end{equation}
%
\begin{equation}
\dot{M}_0 \simeq \left(1\times 10^{11}\;\frac{\text{g}}{\text{sec}}\right)
\left(\frac{M}{M_\odot}\right)^{\!2}
\left(\frac{\rho_\infty}{10^{-24}\;\text{g\,cm}^{-3}}\right)
\left(\frac{T_\infty}{10^4\;\text{K}}\right)^{\!-3/2}\!.
\tag{4.2.15b}
\end{equation}
